\section{Fixed sources, primitive families, and fixed moments}
\label{p1:fixed:section}

The boundary estimate will use large families, but its finitely
many small sizes require stronger arithmetic input. We establish
that input first. All numbers of sources and all moment orders
in this section are fixed before \(N\) tends to infinity.

For a family \(I\) of windows, let \(R(I)\) be the kernel
of the sum map on the effective spaces and put \(\rho_{\mathrm{rel}}(I)=\dim R(I)\).
For \(s\) fixed sources \(\boldsymbol x=(x_1,\ldots,x_s)\) and a
new centre \(z\), write
\begin{equation}
 \delta_z=\dim R(\boldsymbol x,z)-\dim R(\boldsymbol x)
          =\dim\operatorname{pr}_z R(\boldsymbol x,z).
 \label{p1:fixed:delta}
\end{equation}
If the prescribed affine events at the sources and target are
jointly compatible, their conditional probability is
\(p_z2^{\delta_z}\); it is zero otherwise. The same formula applies
to homogeneous prescriptions. The proofs on regular windows and their
effective-space extension in Section~\ref{p13:exact:section} establish
the following inputs without using a distributional approximation.

\subsection{Regularizing the coordinate labels}

The following linear-algebra lemma gives control of every short
circuit of coordinate labels, including repeated labels.

\begin{lemma}[Disjoint recovery]
\label{p1:fixed:recovery}
Let \(R\subseteq\F_2^V\), \(n=|V|\), and let
\(\ell_v\in R^*\) be its coordinate forms. For fixed integers
\(d\ge1\) and \(t\ge2\), there is a subspace \(R'\subseteq R\)
with
\begin{equation}
 \codim_R R'\le 2n t^{-1/d}
 \label{p1:fixed:recovery-cost}
\end{equation}
such that no nonzero \(a\in(R')^*\) has \(t\) representations
\[
 a=\sum_{v\in A_i}\ell_v|_{R'},\qquad
 1\le |A_i|\le d,
\]
by pairwise disjoint subsets \(A_1,\ldots,A_t\) of \(V\).
Each nonzero coordinate label occurs at most \(K=t-1\) times.
For \(1\le j\le d\), a nonzero label has at most
\(j!t^{j-1}\) representations as a sum of \(j\) distinct
linearly independent labels.
\end{lemma}

\begin{proof}
Start with \(W=0\subset R^*\). Whenever \(R^*/W\) has a
nonzero class with \(t\) such disjoint recovery sets, record
those sets and adjoin one representative \(\alpha_i\) of the
class to \(W\). Each step increases \(\dim W\) by one. Fix the
entire resulting sequence of \(m\) steps.

Select each coordinate independently with probability
\(u\in(0,1)\), obtaining \(A\subseteq V\). At step \(i\),
declare success if one of its recorded recovery sets is
contained in \(A\). On success,
\(\alpha_i\in\langle\ell_v:v\in A\rangle+W_{i-1}\).
The failure probability is at most \((1-u^d)^t\), since the
recorded sets at that step are disjoint. For each realization,
\[
 m\le |A|+\#\{\hbox{failed steps}\}.
\]
Taking expectations and choosing \(u=t^{-1/d}\) gives
\[
 m\le \frac{nu}{1-(1-u^d)^t}\le2nt^{-1/d}.
\]
The annihilator \(R'=W^\perp\) has the stated codimension
and satisfies the stopping condition.

The multiplicity bound follows by using singleton recovery sets.
To count independent \(j\)-term representations of a nonzero
\(a\), take a maximal disjoint family of their label supports.
It has fewer than \(t\) members. Its union has at most
\(j(t-1)\) labels and meets every representation. Removing
one of these labels reduces to a \((j-1)\)-term independent
representation of a nonzero vector: the removed label cannot
equal \(a\), because the remaining independent labels would
then sum to zero. Induction from \(j=1\) proves the last bound.
\end{proof}

\subsection{Uniformly few targets of large rank loss}

\begin{proposition}[Large rank loss after fixed sources]
\label{p1:fixed:targets}
For every fixed \(s\ge1\) and \(\eta,\varepsilon>0\),
\begin{equation}
 \sup_{x_1,\ldots,x_s\ {\rm distinct}}
 \#\{z\in I_N\setminus\{x_1,\ldots,x_s\}:
                   \delta_z\ge\eta L\}
       \ll_{s,C,\eta,\varepsilon}N^\varepsilon .
 \label{p1:fixed:targets-bound}
\end{equation}
The sources may be dependent or overlapping.
\end{proposition}

\noindent The complete proof is given in Section~\ref{p14:comp:fixed-targets}.


\begin{proposition}[Conditional sums and propagation]
\label{p1:fixed:propagation}
For every fixed \(s\ge1\) and \(\varepsilon>0\), uniformly
over compatible prescribed affine events at \(s\) distinct
sources,
\begin{equation}
 \sum_{z\notin\{x_1,\ldots,x_s\}}
 \P(J_z=1\mid J_{x_1}=\cdots=J_{x_s}=1)
       \ll_{s,C,\varepsilon}N^\varepsilon .
 \label{p1:fixed:conditional}
\end{equation}
This assertion also holds when the affine prescriptions
are replaced by homogeneous ones. If a family of ordered
\(s\)-tuples has total positive joint mass \(W_{s,N}\),
the mass of its extensions to ordered distinct \(r\)-tuples,
for fixed \(r\ge s\), is
\begin{equation}
 W_{s\to r,N}\ll_{r,C,\varepsilon}N^\varepsilon W_{s,N}.
 \label{p1:fixed:propagate}
\end{equation}
In particular, the mass of fixed \(r\)-tuples containing
a local defect or a compatible contact is
\(O_{r,C,\varepsilon}(N^{-1+\varepsilon})\).
\end{proposition}

\begin{proof}
For \(\delta_z<\eta L\), the conditional sum is at most
\(2Np2^{\eta L}\ll_C N^\eta\). For
\(\delta_z\ge\eta L\), Proposition~\ref{p1:fixed:targets}
counts \(O(N^{\varepsilon/2})\) targets, each of conditional
probability at most one. Take \(\eta=\varepsilon/2\).
The proof uses only the rank loss and compatibility, so
it applies unchanged to homogeneous prescriptions.
Iterating with margin \(\varepsilon/(r-s)\) proves
\eqref{p1:fixed:propagate}. Allowing a seed at any positions
in the ordered tuple multiplies the estimate by a constant
depending on \(r\).

The local and contact seeds have positive masses
\(N^{-1+o(1)}\), by the regular-core estimates in
Section~\ref{p1:arith:section}. Contacts at distance
strictly less than \(L\) have zero joint probability.
Propagate the remaining seeds to obtain the last claim.
\end{proof}

\subsection{The correction to every fixed moment}

\begin{theorem}[Fixed factorial moments]
\label{p1:fixed:moments}
For every fixed integer \(r\ge2\) and every
\(\varepsilon>0\),
\begin{equation}
 \E(Z_N)_r
   =\mu_N^r+\binom r2\mu_N^{r-2}\alpha_N
       +O_{r,C,\varepsilon}(N^{-4/5+\varepsilon}).
 \label{p1:fixed:moment-expansion}
\end{equation}
Equivalently, one may use \(\lambda_N\) and
\(a_N=\E(Z_N)_2-\lambda_N^2\).
All fixed ordinary moments are bounded. More precisely,
with \(T_r(u)=\E[\Pois(u)^r]\),
\begin{equation}
 \E Z_N^r
    =T_r(\mu_N)+\frac{\alpha_N}{2}T_r''(\mu_N)
       +O_{r,C,\varepsilon}(N^{-4/5+\varepsilon}).
 \label{p1:fixed:ordinary}
\end{equation}
For every fixed \(r\ge3\), the \(r\)-th factorial
cumulant is \(O_{r,C,\varepsilon}(N^{-4/5+\varepsilon})\).
\end{theorem}

This first estimate is proved independently of the law approximation.
Its bounded-moment conclusion is used below in the sharp lower bound.
After that reconstruction, Corollary~\ref{p13:optimal:weighted}
improves the scalar moment remainders to
\(O_{r,C,\varepsilon}(N^{-4/3+\varepsilon})\) when the exact
parameters are retained.

\noindent The complete proof is given in Section~\ref{p14:comp:fixed-moments}.


For \(r\ge2\), the first discrepancy in
\eqref{p1:fixed:moment-expansion} has the exact scale
\(L N^{-2/3}\), by \eqref{p1:eq:alpha-leading}.
The theorem fixes \(r\) first. Its proof does not
allow \(r\) to grow like \(\log N/\log\log N\).
The uniform certificate tools in Section~\ref{p14:cert:section} and the
large-order reconstruction in Section~\ref{p13:transfer:section}
supply that separate input.
